\section*{Bab \ref{ch:g12:ph-conductivity-titrations} --- Titrasi dengan pH dan Konduktivitas}

\begin{solution}{exo:g12:ph-conductivity-titrations:1}
$V_E = \qty{20.0}{mL}$; $c = 0.100 \times 20.0 / 20.0 = \qty{0.100}{mol/L}$.
\end{solution}

\begin{solution}{exo:g12:ph-conductivity-titrations:2}
Pada \omterm{def:g12:ph-conductivity-titrations:ph-metric}{setengah ekuivalen} \omterm{def:g9:acids-bases-ph:ph}{pH} $= pK_a$: $pK_a = 3.75$, dekat dengan asam
metanoat (3.74).
\end{solution}

\begin{solution}{exo:g12:ph-conductivity-titrations:3}
Fenolftalein: trayeknya 8.0--10.0 terletak di dalam lonjakan, yang
bersifat basa. Merah metil (4.2--6.3) akan berubah di dekat \omterm{def:g12:ph-conductivity-titrations:ph-metric}{setengah
ekuivalen}, terlalu cepat.
\end{solution}

\begin{solution}{exo:g12:ph-conductivity-titrations:4}
Tanggapan elektrodenya bergeser seiring waktu dan suhu: elektrode itu
diatur agar menunjukkan nilai yang benar dalam dua \omterm{def:g12:buffers-predominance:buffer}{larutan penyangga}
yang \omterm{def:g9:acids-bases-ph:ph}{pH-nya} diketahui.
\end{solution}

\begin{solution}{exo:g12:ph-conductivity-titrations:5}
\ce{H3O+} (349.8) > \ce{OH-} (198.3) > \ce{Cl-} (76.2) > \ce{Na+} (50.3).
\end{solution}

\begin{solution}{exo:g12:ph-conductivity-titrations:6}
$c = 0.050 \times 14.6 / 20.0 = \qty{0.0365}{mol/L}$.
\end{solution}

\begin{solution}{exo:g12:ph-conductivity-titrations:7}
Kemiringan: 2.0, 4.5, 17.5, dan \qty{2.5}{mL^{-1}}. Yang terbesar
terletak antara 10.0 dan \qty{10.2}{mL}: $V_E \approx \qty{10.1}{mL}$.
\end{solution}

\begin{solution}{exo:g12:ph-conductivity-titrations:8}
Sebelum \omterm{def:g11:titration:equivalence}{titik ekuivalen}, setiap \ce{OH-} yang ditambahkan menyingkirkan
sebuah \ce{H3O+} (sangat menghantarkan) dan membawa sebuah \ce{Na+}
(jauh kurang menghantarkan): \omterm{def:g12:ph-conductivity-titrations:conductivity}{konduktivitasnya} turun. Sesudahnya,
\ce{Na+} dan \ce{OH-} menumpuk tanpa bereaksi: \omterm{def:g12:ph-conductivity-titrations:conductivity}{konduktivitasnya} naik.
\end{solution}

\begin{solution}{exo:g12:ph-conductivity-titrations:9}
Di awal, \omterm{def:g9:ions:ion}{ionnya} sedikit (\omterm{def:g12:ka-and-pka:strong-acid}{asam lemah} hampir tidak bereaksi dengan air):
\omterm{def:g12:ph-conductivity-titrations:conductivity}{konduktivitasnya} rendah. Sebelum \omterm{def:g11:titration:equivalence}{titik ekuivalen}, \omterm{def:g9:ions:ion}{ion} etanoat dan \omterm{def:g9:ions:ion}{ion}
natrium terbentuk menggantikan \omterm{def:g7:atoms-and-molecules:molecule}{molekul-molekul} yang tidak bermuatan:
\omterm{def:g12:ph-conductivity-titrations:conductivity}{konduktivitasnya} naik, perlahan. Sesudah \omterm{def:g11:titration:equivalence}{titik ekuivalen}, \ce{Na+} dan
\ce{OH-} menumpuk: \omterm{def:g12:ph-conductivity-titrations:conductivity}{konduktivitasnya} naik lebih cepat. Dua garis yang
naik, dengan garis kedua lebih curam.
\end{solution}

\begin{solution}{exo:g12:ph-conductivity-titrations:10}
Sekitar 2.9, 4.76, dan 8.7. Asam etanoat adalah \omterm{def:g12:ka-and-pka:strong-acid}{asam lemah}: hanya
beberapa persen saja yang menghasilkan \omterm{def:g12:ka-and-pka:oxonium}{ion oksonium}, sehingga \omterm{def:g9:acids-bases-ph:ph}{pH-nya}
dimulai lebih tinggi daripada \omterm{def:g9:acids-bases-ph:ph}{pH} 1.0 asam klorida dengan konsentrasi
yang serupa.
\end{solution}

\begin{solution}{exo:g12:ph-conductivity-titrations:11}
Agar volume \omterm{def:g11:titration:titration}{titran} yang ditambahkan hampir tidak mengubah volume
totalnya: konsentrasinya, dan karena itu \omterm{def:g12:ph-conductivity-titrations:conductivity}{konduktivitasnya}, lalu berubah
mengikuti garis lurus.
\end{solution}

\begin{solution}{exo:g12:ph-conductivity-titrations:12}
\omterm{def:g9:ions:ion}{Ion-ion} hidroksida mula-mula bereaksi dengan \omterm{def:g12:ka-and-pka:strong-acid}{asam kuat} (\ce{H3O+}), lalu
dengan \omterm{def:g12:ka-and-pka:strong-acid}{asam lemah}. Lonjakan pertama menandai habisnya asam klorida
(\omterm{def:g11:titration:equivalence}{volume ekuivalennya} mengukur asam itu), lonjakan kedua menandai
habisnya asam etanoat (volume di antara kedua lonjakan mengukurnya).
\end{solution}

\begin{solution}{exo:g12:ph-conductivity-titrations:13}
Awal: $(349.8 + 76.2) \times 0.0100 = \qty{4.26}{mS/cm}$. Pada \omterm{def:g11:titration:equivalence}{titik
ekuivalen}: $[\ce{Na+}] = [\ce{Cl-}] = 1.00/110 = \qty{0.00909}{mol/L}$,
jadi $(50.3 + 76.2) \times 0.00909 = \qty{1.15}{mS/cm}$. Keduanya
sesuai dengan gambar.
\end{solution}

\begin{solution}{exo:g12:ph-conductivity-titrations:14}
Tidak ada pada $V_E$: lonjakannya terjadi pada volume yang sama, berapa
pun pergeseran pembacaannya. $pK_a$ yang dibaca pada \omterm{def:g12:ph-conductivity-titrations:ph-metric}{setengah ekuivalen}
0.2 terlalu tinggi.
\end{solution}

\begin{solution}{exo:g12:ph-conductivity-titrations:15}
\omterm{def:g9:acids-bases-ph:ph}{pH} tidak berubah selama pengendapan ini, tetapi \omterm{def:g9:ions:ion}{ion-ionnya} berubah:
sebelum \omterm{def:g11:titration:equivalence}{titik ekuivalen}, setiap \ce{Ag+} yang ditambahkan menyingkirkan
sebuah \ce{Cl-} dan membawa sebuah \ce{NO3-} yang \omterm{def:g12:ph-conductivity-titrations:conductivity}{konduktivitasnya}
serupa, sehingga \omterm{def:g12:ph-conductivity-titrations:conductivity}{konduktivitasnya} hampir tetap (\omterm{def:g9:ions:ion}{ion-ion} natrium sudah
ada sejak awal); sesudah \omterm{def:g11:titration:equivalence}{titik ekuivalen}, \ce{Ag+} dan \ce{NO3-}
menumpuk dan \omterm{def:g12:ph-conductivity-titrations:conductivity}{konduktivitasnya} naik. Sebuah garis datar, lalu sebuah
garis yang naik.
\end{solution}

\begin{solution}{pb:g12:ph-conductivity-titrations:1}
\textbf{1.} \qty{6}{g} asam etanoat per \qty{100}{g} cuka.

\textbf{2.} \qty{101}{g} cuka mengandung $6.0 \times 101/100 =
\qty{6.06}{g}$; $M = \qty{60.0}{g/mol}$: \qty{0.101}{mol} dalam
\qty{0.1000}{L}, yaitu \qty{1.01}{mol/L}.

\textbf{3.} Tanpa diencerkan, \qty{10.0}{mL} akan memerlukan sekitar
\qty{101}{mL} natrium hidroksida \qty{0.100}{mol/L}: lebih banyak
daripada isi sebuah buret.

\textbf{4.} Pipet volume \qty{10.0}{mL} dan labu ukur
\qty{100.0}{mL}.

\textbf{5.} \ce{CH3COOH + OH- -> CH3COO- + H2O}.

\textbf{6.} $V_E = \qty{10.1}{mL}$.

\textbf{7.} $c = 0.100 \times 10.1 / 10.0 = \qty{0.101}{mol/L}$.

\textbf{8.} $10 \times 0.101 = \qty{1.01}{mol/L}$.

\textbf{9.} \omterm{def:g9:acids-bases-ph:ph}{pH} 4.76 pada \qty{5.05}{mL}: $pK_a$ asam etanoat.

\textbf{10.} Pada \omterm{def:g11:titration:equivalence}{titik ekuivalen} larutannya mengandung natrium
etanoat, dan etanoat adalah \omterm{def:g12:ka-and-pka:strong-acid}{basa lemah}.

\textbf{11.} Fenolftalein, yang trayeknya 8.0--10.0 mencakup \omterm{def:g9:acids-bases-ph:ph}{pH} pada
\omterm{def:g11:titration:equivalence}{titik ekuivalen} (sekitar 8.7).

\textbf{12.} \omterm{def:g9:ions:ion}{Ion} natrium dan \omterm{def:g9:ions:ion}{ion} etanoat.

\textbf{13.} Setiap \ce{OH-} yang ditambahkan mengubah sebuah \omterm{def:g7:atoms-and-molecules:molecule}{molekul}
yang hampir tidak terionisasi menjadi \ce{CH3COO-}, bersama sebuah
\ce{Na+}: \omterm{def:g9:ions:ion}{ion-ion} yang \omterm{def:g12:ph-conductivity-titrations:conductivity}{konduktivitasnya} sedang ditambahkan secara
perlahan.

\textbf{14.} Sesudah \omterm{def:g11:titration:equivalence}{titik ekuivalen}, \omterm{def:g9:ions:ion}{ion-ion} hidroksida, yang sangat
menghantarkan, menumpuk bersama \omterm{def:g9:ions:ion}{ion-ion} natrium.

\textbf{15.} \omterm{def:g12:ka-and-pka:strong-acid}{Asam lemah} hanya memberikan sangat sedikit \omterm{def:g9:ions:ion}{ion} ke dalam
air.

\textbf{16.} Pada $V_E = \qty{10.1}{mL}$, seperti pada \omterm{def:g12:ph-conductivity-titrations:ph-metric}{titrasi
pH-metri}.

\textbf{17.} $1.01 \times 60.0 = \qty{60.6}{g/L}$.

\textbf{18.} \qty{6.06}{g} dalam \qty{100.0}{mL}.

\textbf{19.} $6.06 \times 100 / 101 = \qty{6.0}{g}$ per \qty{100}{g}.

\textbf{20.} Cuka itu memiliki keasaman \qty{6.0}{\percent}: labelnya
benar.
\end{solution}
